\section{Introduction}

Natural-language prompts now participate directly in the security boundary of
code generation. They usually specify what a program should do, yet leave
validation, trust boundaries, safe APIs, or authorization constraints implicit.
Consequently, code can appear functionally complete while retaining security
weaknesses, and developers can overestimate the security of assisted solutions
\cite{pearce2022copilot,perry2023assistants}. For prompt authors, benchmark
designers, and tool builders, the actionable question is therefore:
\emph{which prompt-side security requirements actually change generated-code
security, for which tasks, and without sacrificing functionality?}

Existing evidence does not answer this question directly. Attribution,
concept-based interpretation, and rationale evaluation associate input text
with model behavior
\cite{sundararajan2017integrated,kim2018tcav,lundberg2017shap,deyoung2020eraser},
while secure-prompt studies compare predefined reminders, reasoning strategies,
and repair loops~\cite{tony2025prompting,bruni2025promptbenchmark}. These views
can expose salient text or useful strategies, but they do not establish that a
specific, context-applicable requirement caused a security change. Conversely,
free-form prompt perturbations can change task meaning, security content, and
presentation together, obscuring what was actually tested.

The missing link is a controlled path from prompt meaning to requirement-level
effect evidence. A natural-language request must first be represented so that
task context and editable security content remain distinct. Observational
evidence can then prioritize a limited number of hypotheses, but cannot be
promoted into an effect claim. Finally, each selected requirement, operation,
policy realization, target outcome, and analysis rule must be frozen before
held-out randomized evaluation.

We introduce \model (\modelfull), which follows this path in three stages.
\textbf{Typed Prompt Representation} maps a programming request into a
\prompttsg and derives context-conditioned requirement templates.
\textbf{Requirement Hypothesis Discovery and Selection} uses TSG-derived
constraints, FCI evidence, and stability under semantic-cluster resampling to
freeze a budgeted hypothesis set. \textbf{Evaluation of Requirement Effects}
constructs matched prompt policies and estimates their assigned-arm effects on
cluster-disjoint held-out tasks. The boundary is deliberate: the Prompt TSG
represents candidate requirements; observational discovery determines which
hypotheses to test; randomized assignment determines the effect of the frozen
prompt policy.

We use \emph{matter} as a staged, evidence-bound conclusion rather than a loose
synonym for importance. A requirement is \emph{effect-relevant} when the
multiplicity-adjusted target-versus-no-op assigned-arm interval excludes zero
within its frozen scope. It is \emph{beneficial} when that supported effect has
the security-improving direction---for example, ADD improves secure-code yield
or REMOVE reduces it. It is \emph{practically actionable} only when the
beneficial effect is also target-specific and functionality-preserving under
the frozen decision rules. Inconclusive evidence is not evidence of no effect,
and none of these labels generalizes beyond the evaluated context, operation,
realization distribution, generator model, and outcome.

This paper makes four contributions:

\begin{itemize}
  \item \textbf{A requirement-centered prompt representation.} A typed Prompt
  TSG separates immutable task-security context from editable safety factors
  and derives applicability-aware ADD and REMOVE templates from a finite
  catalog.
  \item \textbf{A budgeted discovery-to-freeze procedure.} Typed constraints,
  FCI uncertainty, and semantic-cluster stability prioritize requirement
  hypotheses while keeping observational selection evidence distinct from
  randomized effect evidence.
  \item \textbf{A requirement-level randomized evaluation.} Matched prompt
  policies, blocked assignment, and semantic-task-clustered ITT expose effect,
  specificity, and functionality as separate axes, retaining null-compatible,
  opposite, failed, unknown, and non-evaluable outcomes required by the frozen
  protocol.
  \item \textbf{A usable evidence interface.} A standardized requirement
  evidence record joins context, operation, selection trace, randomized effect,
  uncertainty, and scope; a separately governed study evaluates experts'
  perceived utility without treating those ratings as effect validation.
\end{itemize}
